Understanding which properties of benchmark datasets predict ML reproducibility failure could enable automated pre-review auditing of high-risk papers --- but only if the required dataset metadata is accessible at scale. We investigate whether HuggingFace Hub dataset cards and OpenML run counts can serve as data sources for a metadata-based reproducibility auditing study on Raff's 255-paper corpus of top-venue ML papers from 1984--2017. Our infrastructure feasibility study reveals a systematic coverage gap: HF Hub achieves only 30\% dataset card coverage for the corpus (concentrated in NLP benchmarks and small-scale image classification benchmarks, with large-scale DL vision, speech, and graph datasets absent), while OpenML provides pre-publication entries for only 22\% (restricted to classical tabular datasets). Neither threshold is sufficient for valid regression analysis, and both gaps trace to the same cause --- each platform's ecosystem reflects its founding community rather than the broader benchmark landscape of pre-2018 ML literature. These findings show that any reproducibility auditing study of pre-2018 ML papers built on HF+OpenML will be structurally biased toward NLP papers and will miss the deep learning benchmarks that dominate the corpus. We characterize the root causes, provide a validated data acquisition pipeline for the community, and identify Papers With Code as the appropriate primary data source for a redesigned study --- one that the theoretical hypothesis linking dataset misuse signals to reproducibility failure still awaits and merits.
